% Extracellular depletion: fixed removal --- plan \S3.3, "\S4".
% Sources: 03_fixed_removal.tex:9 (moved here as the opening), :23-80,
% :86-95, :99-125, :132-151, :157-158, :161-199, :201-214, :217-228,
% :232-257.
% Renames applied: theta -> \vartheta, R_0 -> \mathcal{R}(\vartheta),
% \mathcal{K} -> \mathcal{M}, \mathcal{G} -> F(\vartheta), I_- -> b,
% V_{inf} -> V_\infty, E(Y_t) -> E(W_t).
% Corrections applied (plan \S5.1): (a) the smaller root at 03:205;
% (b) the summation limit at 03:214, replaced by the exact complement;
% (d) the two missing summands at 03:223-224; (e) Pr(r>=0) at 03:171.
% Both \cancel constructions deleted.
%
% PRESERVED NEAR-VERBATIM, do not restyle (plan \S3.6, report item 2):
%   the leading-edge passage at 03:88.

\section{Extracellular depletion: fixed removal}
\label{path:sec:removal}

The second nonlinearity is extracellular, and it enters as an extension to
the chained-rupture model of \cref{path:sec:chainedruptures}. The extension
is one assumption, arising from the same argument of immune depletion: that
following a burst release event, some defined fixed count of the agents is
removed. From that assumption alone, the question of whether bursting
strategies might be of fundamental advantage can be put to a calculation.

\subsection{The playing field}
\label{path:sec:playingfield}

Consider again the concept of spatial drift, which for all its imperfection
improves on the simpler dynamic of chained immediate transfer. Picture the
playing field: a three-dimensional extracellular matrix populated by three
types of agent --- replicative units, that is pathogen; host cells, here
macrophages; and another type of immune-active cell, which we shall call
simply \emph{killer cells}. This third type is able to eliminate units of
pathogen on contact, being used up in the process: one-to-one encounters in
which a killer cell removes a replicative agent and itself from play.

To begin with, the extracellular matrix is populated evenly with target
cells and killer cells, each group occupying the space at a uniform spatial
density. Killer cells are smaller and more numerous than target cells, a
greater number per unit volume. Dynamics within target cells are as
understood in the birth--death--catastrophe process: the units of pathogen
replicate and die within the cell until eventually rupture is caused and the
load is released.

In this model, movement is taken to be Brownian. An extracellular colony
will drift from its burst location, spreading out and coming into contact
with surrounding cells. Some will encounter target cells, others killer
cells. The intention here is not to construct a full mathematical
description of this game, which would be cumbersome and intractable; it can
be simulated computationally, and that simulation is what supports the
assumption the model is used to justify. Suppose that no new killer cells
are introduced into the environment, so that pathogen movement and
absorption occur on a shorter timescale than killer generation. Then, with
greater burst sizes, each individual bacterium will have a smaller
probability of being removed. Those at the leading edge, as it were,
clearing the way for their comrades in the rear. \Cref{path:fig:loads} puts
a large and a small release of that simulation side by side.

\begin{figure}[htbp]
    \centering
    \setlength{\tabcolsep}{1.5pt}
    \begin{tabular}{@{}cccc@{}}
      & \multicolumn{3}{c}{\footnotesize time $\longrightarrow$} \\[1pt]
      \rotatebox{90}{\footnotesize\itshape large load}
      & \includegraphics[width=0.295\textwidth]{big1.png}
      & \includegraphics[width=0.295\textwidth]{big2.png}
      & \includegraphics[width=0.295\textwidth]{big3.png} \\
      & {\footnotesize (a)} & {\footnotesize (b)} & {\footnotesize (c)} \\[3pt]
      \rotatebox{90}{\footnotesize\itshape small load}
      & \includegraphics[width=0.295\textwidth]{small1.png}
      & \includegraphics[width=0.295\textwidth]{small2.png}
      & \includegraphics[width=0.295\textwidth]{small3.png} \\
      & {\footnotesize (d)} & {\footnotesize (e)} & {\footnotesize (f)} \\
    \end{tabular}
    \caption{Two runs of the same simulation differing in released number
    alone; every rate is the same. Depletable extracellular immune factors
    are consumed by the agents they remove, so a large load creates a safety
    in numbers --- a ``flooding effect'' --- and the killers that act early
    are not available later. The comparison to draw is between~(c) and~(f).
    In~(c) the large load has cleared a patch of the killer field and
    survivors remain inside it, poised on the surrounding host cells; in~(f)
    the small load is exhausted before the field is, and nothing persists to
    the next round of absorption and replication.}
    \label{path:fig:loads}
\end{figure}

If free agents do tend to find a new host before long --- an assumption
taken to its limit in the case of chained immediate transfer --- then this
observation offers the basis of an approximate mathematical rule. Burst size
and the number of pathogen killed in the process of transfer stand in some
relation to one another, and across larger and larger bursts the casualty
count should rise ever more slowly.
Qualitatively this might be a logarithmic progression,
\begin{equation}
\vartheta(r)\;\sim\;\log(r+1),
\label{path:eq:logremoval}
\end{equation}
where $r$ is a given burst size and $\vartheta$ the number of replicative
agents removed by killer cells; the $+1$ simply makes the curve pass
through the origin.

\Cref{path:fig:logremoval} draws that curve, and its flatness is what the
rest of the section idealises. It is reasonable to impose a \emph{fixed
removal} following release, in which the number removed by killer cells is a
specified constant regardless of burst size, except that if the burst size
is smaller than that constant, all of the units are killed. Taking
$\vartheta$ now to be that constant,
\begin{equation}
\text{removal} =
\begin{cases}
r, & r\le\vartheta,\\
\vartheta, & r>\vartheta,
\end{cases}
\qquad\text{that is, } \min(r,\vartheta).
\label{path:eq:fixedremoval}
\end{equation}

\begin{figure}[htbp]
\centering
\begin{subfigure}{0.48\textwidth}
\centering
\begin{tikzpicture}
\begin{axis}[
  width=\textwidth, height=0.66\textwidth,
  axis lines=left,
  xlabel={burst size $r$},
  ylabel={casualties $\vartheta(r)$},
  xmin=0, xmax=50, ymin=0, ymax=4.3,
  xtick={0,10,20,30,40,50},
  ytick={0,1,2,3,4},
  tick label style={font=\footnotesize},
  label style={font=\footnotesize},
]
\addplot[na1, very thick, domain=0:50, samples=200]{ln(x+1)};
\end{axis}
\end{tikzpicture}
\caption{the saturating law \cref{path:eq:logremoval}}
\label{path:fig:logremoval}
\end{subfigure}
\hfill
\begin{subfigure}{0.48\textwidth}
\centering
\begin{tikzpicture}
\begin{axis}[
  width=\textwidth, height=0.66\textwidth,
  axis lines=left,
  xlabel={burst size $r$},
  ylabel={removed},
  xmin=0, xmax=20, ymin=0, ymax=5,
  xtick={0,5,10,15,20},
  ytick={3.5}, yticklabels={$\vartheta$},
  tick label style={font=\footnotesize},
  label style={font=\footnotesize},
]
\addplot[Rup2, dashed, thick, domain=0:20, samples=2]{3.5};
\addplot[na1, very thick, domain=0:20, samples=400]{min(x,3.5)};
\end{axis}
\end{tikzpicture}
\caption{the step idealisation \cref{path:eq:fixedremoval}}
\label{path:fig:capped}
\end{subfigure}
\caption{The casualty count against burst size: the law the biology
suggests, and the caricature the calculation uses. In~(a) the vertical axis
is the number of agents killed, not a rate --- what falls as the burst grows
is the killed \emph{fraction}, since the curve flattens while the burst size
does not. Doubling a load of ten adds well under one casualty, which is the
whole of the flooding intuition in one line. In~(b) that flatness is pushed
to its limit: below the budget every agent dies, and above it the plateau at
$\vartheta$ means a burst pays the same toll however large it is, so all the
marginal yield survives. The reason for adopting the caricature is that it
makes the replication ratio computable in closed form.}
\label{path:fig:removal}
\end{figure}

\Cref{path:fig:capped} shows the idealisation against which the rest of the
section works. Its efficacy rests on the expectation that the same
qualitative dynamics are shared with logarithmic removal, or with whatever
the actual relation turns out to be, provided it is of this sort. Two
further assumptions are needed before the calculation can proceed, and both
are restrictions rather than conveniences.

\begin{remark}[Killer prevalence resets]
\label{path:rem:reset}
At the time of each rupture event the killer cell prevalence is assumed to
have returned to its steady state and to be once again of uniform spatial
distribution. This is what makes $\vartheta$ the same at every delivery. It
requires the interval between ruptures to be long relative to the recovery
of the effector population, whose turnover in an established infection is
itself a modelled quantity \cite{ribeiro2002vdt}; where it is not, the
budget paid at successive deliveries would be correlated, and the
independence used below would fail.
\end{remark}

\begin{remark}[One target cell, or many]
\label{path:rem:onetarget}
The fixed-removal dynamic may be framed in the context of chained transfers,
where all agents released have the same target cell to go to; or in a case
where each successful pathogenic unit goes on to found an infection in a new
target cell. Under the first framing the replication ratio would be the
expected number of surviving released units \emph{per incident agent}. The
second framing is taken here, so that the ratio counts newly infected target
cells, which is what makes it comparable with the reproduction numbers of
\ChPop{}.
\end{remark}

\subsection{The replication ratio}
\label{path:sec:ratio}

Take, then, the second option: each released agent goes on to found a new
infection in a new target cell. The quantity of interest is the effective
\emph{replication ratio} between target cell infections --- the expected
number of newly infected target cells created by a host founded from one
initial bacterium --- and because each surviving released agent infects a new
host, that ratio is simply the expected number of survivors.

Define $r$ as the outcome of a birth--death--catastrophe process with a
single incident unit. We must choose whether to include outcomes that are
zero through internal death, and it makes practical sense here not to ignore
the cases $r=0$; with that choice
\begin{equation}
\ex{r}=\lim_{t\to\infty}\ex{W_t}=V_\infty,
\label{path:eq:Er}
\end{equation}
the lifetime yield of \cref{bdc:eq:Vinf}. The letter $r$ is that of
\ChDist{} as well: its $r_k$ is the load at the $k$th rupture of a chain
(\cref{dist:eq:rklaw}), and the $r$ here is that chain's first term.

The number surviving fixed removal is $\max(0,r-\vartheta)$, so the expected
replication ratio is
\begin{equation}
\mathcal{R}(\vartheta) = \ex{\max(0,\,r-\vartheta)}.
\label{path:eq:Rdef}
\end{equation}
This is not the $R_0$ of \ChPop{}, and \cref{path:sec:settled} states the
relation between them; the change of letter is deliberate.

Conditioning on whether the burst size exceeds the budget gives
\begin{equation}
\ex{\max(0,r-\vartheta)}
= \bigl[\ex{r\mid r\ge\vartheta}-\vartheta\bigr]\pr{r\ge\vartheta},
\label{path:eq:Rcond}
\end{equation}
the term for $r<\vartheta$ vanishing because the surviving count is then
zero. Extracting $\ex{r\mid r\ge\vartheta}$ from the decomposition
\[
\ex{r}=V_\infty
= \ex{r\mid r\ge\vartheta}\pr{r\ge\vartheta}
+ \ex{r\mid r<\vartheta}\pr{r<\vartheta}
\]
and cancelling $\pr{r<\vartheta}$ against the denominator inside
$\ex{r\mid r<\vartheta}$ turns \cref{path:eq:Rcond} into a difference of two
partial sums. This is analytically computable because the burst-size law is
known: unconditionally, including the mass $b$ of realisations that never
burst,
\begin{equation}
\pr{r=i}=\frac{\delta}{\lambda}\,a^{-i}\quad (i\ge1),
\qquad
\pr{r=0}=b,
\label{path:eq:burstlaw}
\end{equation}
which is \cref{dist:eq:burstlaw}. Hence
\begin{equation}
\ex{\max(0,r-\vartheta)}
= V_\infty-\mathcal{M}-\vartheta\bigl(1-F(\vartheta)\bigr),
\label{path:eq:Rpartial}
\end{equation}
where
\begin{equation}
F(\vartheta)=\pr{r<\vartheta}=b+\sum_{i=1}^{\vartheta-1}\frac{\delta}{\lambda}a^{-i},
\qquad
\mathcal{M}=\sum_{i=1}^{\vartheta-1} i\,\frac{\delta}{\lambda}a^{-i}.
\label{path:eq:FM}
\end{equation}
Because the formula $\pr{r=i}=(\delta/\lambda)a^{-i}$ does not apply at
$i=0$, the sum for $F$ begins at $1$ and carries the separate term $b$ in
front of it: that is the probability that the process ends in internal
death, an output of zero.

\subsubsection{The law is normalised, and the tail is free}
\label{path:sec:normalisation}

Before analysing \cref{path:eq:Rpartial} it is worth confirming that
\cref{path:eq:burstlaw} is a probability distribution, since the two pieces
come from different places. The probability that $r=0$ is $b$, the smaller
root of the characteristic quadratic,
\begin{equation}
b=\eta-\sqrt{\eta^{2}-\frac{\mu}{\lambda}},
\qquad
\eta=\frac{\lambda+\mu+\delta}{2\lambda},
\label{path:eq:broot}
\end{equation}
as fixed by \cref{bdc:lem:roots}, where $b<1<a$, $ab=\mu/\lambda$ and
$a+b=(\lambda+\mu+\delta)/\lambda$. It remains to show that
$b+\sum_{i\ge1}\pr{r=i}=1$, that is that
\[
\sum_{i=1}^{\infty}\frac{\delta}{\lambda}\lrb{\frac1a}^{i}=1-b .
\]
Using $\sum_{k\ge1}y^{k}=y/(1-y)$ for $|y|<1$, the left-hand side is
$\frac{\delta}{\lambda}\cdot\frac{1/a}{1-1/a}$, so the claim is equivalent to
\[
\frac{\delta}{\lambda}=(1-b)(a-1)=(a+b)-ab-1,
\]
and substituting the two Vieta identities gives
\[
(a+b)-ab-1=\frac{\lambda+\mu+\delta}{\lambda}-\frac{\mu}{\lambda}-1
=\frac{\delta}{\lambda},
\]
as required. The same identity delivers the upper tail with no further work:
\begin{equation}
1-F(\vartheta)=\pr{r\ge\vartheta}=(1-b)\,a^{-(\vartheta-1)},
\label{path:eq:tail}
\end{equation}
which is exact for every integer $\vartheta\ge1$. The tail of the burst law
is geometric, and \cref{path:eq:tail} is the fact from which the next
subsection's closed form follows.

\subsection{Analysis of proliferability}
\label{path:sec:proliferability}

The question is whether $\mathcal{R}(\vartheta)$ is greater or less than
one, and if greater, by how much; and how that changes with $\delta$. Do
larger burst sizes make an infection more viable? The geometric tail of
\cref{path:eq:tail} already answers: a budget spent against a geometric
law removes a fixed proportion of it rather than a fixed amount.

\begin{proposition}[Fixed-removal replication ratio]
\label{path:prop:ratio}
For an integer removal budget $\vartheta\ge0$,
\begin{equation}
\mathcal{R}(\vartheta)=\ex{\max(0,\,r-\vartheta)}=V_\infty\,a^{-\vartheta}.
\label{path:eq:R0}
\end{equation}
Each unit of removal costs the burst a factor $1/a$, independently of
$\vartheta$, of $\mu$, and of everything about the burst-size law except
$a$.
\end{proposition}

\begin{proof}
The burst law \cref{path:eq:burstlaw} is geometric on $\{1,2,\dots\}$ with
ratio $1/a$, hence memoryless above any integer threshold:
\[
\sum_{i>\vartheta}(i-\vartheta)\frac{\delta}{\lambda}a^{-i}
=\frac{\delta}{\lambda}a^{-\vartheta}\sum_{k\ge1}k\,a^{-k}
=a^{-\vartheta}\,\ex{r}
=V_\infty a^{-\vartheta}. \qedhere
\]
\end{proof}

Two consequences follow in one line each, and one restriction.

\begin{corollary}[Threshold]
\label{path:cor:threshold}
$\mathcal{R}(\vartheta)>1$ if and only if
$\vartheta<\vartheta^{\ast}:=\log V_\infty/\log a$. At
$(\lambda,\mu,\delta)=(1,0.2,0.05)$ this gives $\vartheta^{\ast}=44.13$; at
$(1,0.5,1/3)$ it gives $\vartheta^{\ast}=1.71$.
\end{corollary}

\begin{corollary}[Why removal creates a threshold at all]
\label{path:cor:whythreshold}
$V_\infty>1$ for every admissible $(\lambda,\mu,\delta)$, so
$\mathcal{R}(0)>1$ always. The subcritical regime exists only because of the
removal.
\end{corollary}

\begin{remark}[Partial-sum expansion]
\label{path:rem:partial}
The same quantity expanded without memorylessness is
\begin{equation}
\ex{\max(0,r-\vartheta)}
= V_\infty-\vartheta(1-b)
+ \frac{\delta}{\lambda}
  \lrb{\frac{\frac1a\lrb{\vartheta-1-\vartheta\frac1a+\lrb{\frac1a}^{\vartheta}}}
            {\lrb{1-\frac1a}^{2}}} ,
\label{path:eq:Rintermediate}
\end{equation}
which follows from closing the sums in \cref{path:eq:FM} by the identities
$\sum_{i=1}^{n-1}s^{i}=s(1-s^{n-1})/(1-s)$ and
$\sum_{i=1}^{n-1}i\,s^{i}=s(1-ns^{n-1}+(n-1)s^{n})/(1-s)^{2}$. It is exact
at integer $\vartheta$ and is recorded because
\cref{path:app:verification} checks it independently of
\cref{path:eq:R0}.
\end{remark}

\begin{remark}[Scope: $\vartheta$ is an integer]
\label{path:rem:integer}
The derivation sums to $\vartheta-1$ and so requires
$\vartheta\in\mathbb{Z}_{\ge0}$. That is the honest reading in any case,
since $\vartheta$ counts killed particles, and it is adopted as part of the
model definition. For a real budget the corresponding quantity is
$\ex{\max(0,r-\vartheta)}
=a^{-\lceil\vartheta\rceil}\bigl(V_\infty+(\lceil\vartheta\rceil-\vartheta)L\bigr)$
with $L=a(1-b)$, which reduces to \cref{path:eq:R0} exactly when $\vartheta$
is an integer; \cref{path:app:verification} records its accuracy.
\end{remark}

The behaviour of \cref{path:eq:R0} in $\delta$ is drawn in
\cref{path:fig:ratio}.
The reading that matters for what follows is the geometry of the crossing:
because $\mathcal{R}(\vartheta)=V_\infty a^{-\vartheta}$ is a product of a
term that grows as $\delta$ falls and a penalty that is exponential in the
budget, an increase in the toll can always be answered by a large enough
burst, and the exchange rate between the two is exactly one factor of $a$
per unit of $\vartheta$.

\begin{figure}[htbp]
\centering
\includegraphics[width=0.98\textwidth]{ratio_removal.pdf}
\caption{The replication ratio \cref{path:eq:R0} against $\delta$ at
$\lambda=0.8$, $\mu=0$, for budgets $\vartheta=0$, $1$ and $2$, with
$V_\infty$ dashed for reference throughout. In panel~(a) the two curves
coincide, since at zero removal the surviving release is the whole release;
that panel has no crossing, which is \cref{path:cor:whythreshold} drawn.
Elsewhere the marked point is the crossing $\mathcal{R}=1$ of
\cref{path:cor:threshold}, at $\delta=0.800$ and $\delta=0.494$: to the left
of it the infection proliferates, to the right the removal extinguishes it.
Each further unit of budget moves the crossing left, so a pathogen facing a
heavier toll needs a lower catastrophe rate --- a larger burst --- to stay
supercritical.}
\label{path:fig:ratio}
\end{figure}

\FloatBarrier
